% Modèle de rapport d'activité — évaluation à mi-vague (CNRS, Comité national).
% Plan et consignes : « Recommandations pour la rédaction du dossier d'évaluation de l'activité des chercheurs et des chercheuses du CNRS », SGCN / DRH, novembre 2025.
% Remplacer les \consigne{…} par votre texte, puis compiler avec pdflatex (ou Overleaf).
\documentclass[10pt,a4paper]{article}
\usepackage[utf8]{inputenc}
\usepackage[T1]{fontenc}
\usepackage[french]{babel}
\frenchsetup{StandardLists=true}
\usepackage[a4paper,margin=2cm]{geometry}
\usepackage{helvet}            % Arial ou équivalent, comme demandé
\renewcommand{\familydefault}{\sfdefault}
\usepackage{xcolor,enumitem,titlesec,hyperref}
\definecolor{cnrs}{HTML}{22397F}
\definecolor{consigne}{HTML}{6B7280}
\setlength{\parindent}{0pt}
\setlength{\parskip}{0.4em}
\setlist{topsep=2pt,itemsep=1pt}
\titleformat{\section}{\large\bfseries\color{cnrs}}{\thesection.}{0.5em}{}
\titleformat{\subsection}{\normalsize\bfseries\color{cnrs}}{}{0em}{}
\newcommand{\consigne}[1]{{\color{consigne}\itshape #1}}

\begin{document}

{\LARGE\bfseries\color{cnrs} Rapport d'activité --- évaluation à mi-vague}\par\medskip
\textbf{Nom, prénom :} …\par
\textbf{Grade :} …\par
\textbf{Unité (UMR) :} …\par
\textbf{Section(s) / CID :} …\par
\textbf{Période de référence :} …\par\medskip
\consigne{Période de référence : les 5 derniers semestres, indiquée dans le courriel d'ouverture de la campagne d'évaluation. Mise en forme : Arial 10, interligne 1. Conserver l'ordre des rubriques, même si certaines restent vides. Dépôt en PDF dans l'application e-valuation.}

\section{Curriculum vitæ}
\consigne{Votre curriculum vitæ.}

\section{Résumé de votre activité}
\consigne{Résumé (1/2 page) qui permettra à l'ensemble des membres de la Section et/ou CID de prendre connaissance de votre profil scientifique.}

\section{Rapport d'activité}
\consigne{Vous présenterez de manière critique, en 15 pages maximum, les éléments suivants, en veillant à faire ressortir la valeur et l'influence de votre travail sur l'avancée des connaissances dans votre domaine :}
\begin{itemize}
  \item \consigne{votre activité de recherche au cours des 5 derniers semestres, les facteurs qui ont guidé vos choix, les résultats obtenus, leur portée et leur impact, et les éventuelles difficultés rencontrées ;}
  \item \consigne{votre participation à des collaborations avec des partenaires académiques, contractuelles ou non, dans un cadre français, européen ou international (préciser votre rôle, les partenaires, les thèmes des travaux et leur impact) ;}
  \item \consigne{vos participations à des contrats de recherche (ANR, contrats européens, autres agences) : intitulé long, acronyme, durée, partenaires et montants, nature de vos responsabilités (PI, co-PI, Work Package Manager), thème et contenu des travaux, portée et impact ;}
  \item \consigne{la place de votre recherche dans celle de votre unité, du point de vue des thématiques et des équipes (éventuellement sous forme d'organigramme) ;}
  \item \consigne{vos mobilités éventuelles, géographiques, thématiques ou fonctionnelles (séjours dans d'autres laboratoires publics ou privés, en France, en Europe ou à l'étranger, et leurs apports) ;}
  \item \consigne{les distinctions scientifiques reçues et les contributions qu'elles ont récompensées.}
\end{itemize}
\subsection*{Productions scientifiques (5 au maximum)}
\consigne{Classez vos productions par ordre d'importance. Pour chacune, expliquez sa portée, son impact et votre contribution personnelle, et donnez le lien vers la production elle-même dans HAL (ou une autre archive ouverte).}
\begin{enumerate}
  \item \consigne{Référence complète — lien HAL}\\ \consigne{Portée, impact et contribution personnelle : …}
  \item \consigne{Référence complète — lien HAL}\\ \consigne{Portée, impact et contribution personnelle : …}
\end{enumerate}

\section{Enseignement, formation et diffusion de la culture scientifique}
\consigne{Pour les 5 derniers semestres :}
\begin{itemize}
  \item \consigne{thèses dirigées et post-doctorats encadrés (noms, sujets, part prise dans l'encadrement, contenu et portée des travaux, devenir des docteurs et post-doctorants) ;}
  \item \consigne{participations à l'enseignement (niveau, titres des cours, établissements, nombre annuel d'heures) ;}
  \item \consigne{implication dans le partage, le management ou la publication des données de la recherche ;}
  \item \consigne{organisation de conférences, workshops, congrès et rencontres « portes ouvertes » ;}
  \item \consigne{revues ou ouvrages de vulgarisation et diffusion de l'information scientifique vers des publics non spécialistes ;}
  \item \consigne{interventions dans la presse écrite et audiovisuelle ;}
  \item \consigne{travaux d'expertise ;}
  \item \consigne{autres.}
\end{itemize}

\section{Transfert technologique, relations industrielles et valorisation}
\consigne{Pour les 5 derniers semestres :}
\begin{itemize}
  \item \consigne{contrats industriels et autres (durée, financeurs, partenaires et montants, nature de vos responsabilités, thème, contenu, portée et impact des travaux) ;}
  \item \consigne{projets de création d'entreprises ;}
  \item \consigne{travaux donnant lieu à des brevets ou à des développements de logiciels (votre rôle, contenu, difficultés rencontrées, impact) ;}
  \item \consigne{activités de consultance (destinataires, temps passé) ; distinctions reçues pour des activités de transfert ou de valorisation ;}
  \item \consigne{autres.}
\end{itemize}

\section{Encadrement, animation et management de la recherche}
\consigne{Pour les 5 derniers semestres, en mettant en lumière la manière d'exercer vos responsabilités, vos constats et les difficultés éventuelles rencontrées :}
\begin{itemize}
  \item \consigne{animation de programmes ou projets français, européens ou internationaux (type, objectif, ampleur, impact) ;}
  \item \consigne{direction d'équipe ou de laboratoire (nombre de personnes) ;}
  \item \consigne{autres responsabilités ou activités collectives (management de la recherche, missions d'intérêt général, conseils scientifiques, commissions, comités de lecture, etc.) ;}
  \item \consigne{autres.}
\end{itemize}

\section{Objectifs / Projet de recherche}
\consigne{Sans objet pour une évaluation à mi-vague (rubrique conservée pour respecter le plan).}

\section{Observations éventuelles}
\consigne{Observations que vous souhaitez porter à la connaissance de votre Section et/ou CID (1/2 page).}

\end{document}
